\documentclass[12pt]{article}
\usepackage[margin=1in]{geometry}
\usepackage{amsmath}
\usepackage{amssymb}
\usepackage{enumitem}
\usepackage{xcolor}
\usepackage{tcolorbox}

\title{The Independence Axiom in Expected Utility Theory}
\author{ECO310 Midterm Review}
\date{March 5, 2026}

\begin{document}

\maketitle

\section{Introduction to Expected Utility}

Expected Utility Theory provides a framework for understanding decision-making under uncertainty. For preferences over lotteries to be represented by expected utility, they must satisfy four key conditions:

\begin{enumerate}
    \item \textbf{Completeness}: For any two lotteries $p$ and $q$, either $p \succeq q$ or $q \succeq p$ (or both)
    \item \textbf{Transitivity}: If $p \succeq q$ and $q \succeq r$, then $p \succeq r$
    \item \textbf{Continuity}: Small changes in probabilities lead to small changes in preferences
    \item \textbf{Independence Axiom}: The critical condition for expected utility
\end{enumerate}

\section{The Independence Axiom}

\subsection{Formal Statement}

The independence axiom states that if you prefer lottery $p$ to lottery $q$, then you should still prefer $p$ to $q$ even when both are mixed with some third lottery $r$ in the same proportion.

\begin{tcolorbox}[colback=blue!5!white,colframe=blue!75!black,title=Independence Axiom]
If $p \succ q$, then for any lottery $r$ and probability $\alpha \in (0,1)$:
$$\alpha p + (1-\alpha)r \succ \alpha q + (1-\alpha)r$$
\end{tcolorbox}

\subsection{Intuition}

The independence axiom captures the idea that your preference between two lotteries should not depend on what happens in states that are common to both. If you prefer $p$ to $q$, mixing both with the same lottery $r$ shouldn't reverse your preference.

Think of it this way: if you prefer apples to oranges, then you should still prefer "80\% chance of apples, 20\% chance of bananas" to "80\% chance of oranges, 20\% chance of bananas."

\section{Sally's Example: Violating Independence}

\subsection{The Setup}

Consider three outcomes with monetary values:
\begin{itemize}
    \item $x_1 = \$1,000,000$
    \item $x_2 = \$5,000,000$
    \item $x_3 = \$0$
\end{itemize}

Sally expresses the following preferences over lotteries (where each lottery is written as $(P(x_1), P(x_2), P(x_3))$):

\begin{align}
\text{Preference 1:} \quad &(0, 1, 0) \succ (0.01, 0.89, 0.10) \label{eq:pref1}\\
\text{Preference 2:} \quad &(0.90, 0, 0.10) \succ (0.89, 0.11, 0) \label{eq:pref2}
\end{align}

\subsection{Interpreting the Preferences}

\textbf{Preference 1:} Sally prefers a guaranteed \$5,000,000 to a lottery with:
\begin{itemize}
    \item 1\% chance of \$1,000,000
    \item 89\% chance of \$5,000,000
    \item 10\% chance of \$0
\end{itemize}

This shows Sally is somewhat risk-averse—she'd rather have the sure thing than take a 10\% chance of getting nothing.

\textbf{Preference 2:} Sally prefers a lottery with 90\% chance of \$1,000,000 and 10\% chance of \$0 to a lottery with 89\% chance of \$1,000,000 and 11\% chance of \$5,000,000.

This seems to show Sally is willing to give up a potential extra \$4,000,000 to increase her chance of winning something from 89\% to 90\%.

\subsection{The Key Insight}

Here's the crucial observation: \textbf{the two left-hand-side lotteries add up to the same total as the two right-hand-side lotteries}.

Let's verify:
\begin{align*}
(0, 1, 0) + (0.90, 0, 0.10) &= (0.90, 1.00, 0.10)\\
(0.01, 0.89, 0.10) + (0.89, 0.11, 0) &= (0.90, 1.00, 0.10)
\end{align*}

They're identical!

\section{Proving the Violation}

\subsection{What Expected Utility Would Require}

If Sally's preferences could be represented by expected utility, there would exist a utility function $u$ such that preferences are determined by expected utility values.

\textbf{From Preference 1:}
\begin{align*}
u(x_2) &> 0.01u(x_1) + 0.89u(x_2) + 0.10u(x_3)\\
u(x_2) - 0.89u(x_2) &> 0.01u(x_1) + 0.10u(x_3)\\
0.11u(x_2) &> 0.01u(x_1) + 0.10u(x_3)
\end{align*}

\textbf{From Preference 2:}
\begin{align*}
0.90u(x_1) + 0.10u(x_3) &> 0.89u(x_1) + 0.11u(x_2)\\
0.90u(x_1) - 0.89u(x_1) + 0.10u(x_3) &> 0.11u(x_2)\\
0.01u(x_1) + 0.10u(x_3) &> 0.11u(x_2)
\end{align*}

\subsection{The Contradiction}

We have derived:
\begin{align*}
\text{From Preference 1:} \quad &0.11u(x_2) > 0.01u(x_1) + 0.10u(x_3)\\
\text{From Preference 2:} \quad &0.01u(x_1) + 0.10u(x_3) > 0.11u(x_2)
\end{align*}

\begin{tcolorbox}[colback=red!5!white,colframe=red!75!black,title=Contradiction!]
These two inequalities cannot both be true! We have:
$$0.11u(x_2) > 0.01u(x_1) + 0.10u(x_3) > 0.11u(x_2)$$
which implies $0.11u(x_2) > 0.11u(x_2)$, a logical impossibility.
\end{tcolorbox}

\subsection{Why Independence Fails}

The independence axiom would say: if we can write Preference 1 as a mixture of two lotteries where one lottery appears in Preference 2, then the preferences should be consistent.

Notice that we can decompose the lotteries as:
\begin{align*}
(0, 1, 0) &= 0.11 \times (0, 1, 0) + 0.89 \times (0, 1, 0)\\
(0.01, 0.89, 0.10) &= 0.11 \times (0, 0, 1) + 0.89 \times (0, 1, 0) + 0.01 \times (1, 0, 0)
\end{align*}

The mixing structure reveals the inconsistency when combined with Preference 2.

\section{Implications}

\subsection{Real-World Behavior}

Sally's preferences are an example of the famous \textbf{Allais Paradox}, which demonstrates that real people often violate the independence axiom. This suggests:

\begin{itemize}
    \item Expected utility theory may not perfectly describe actual human decision-making
    \item People treat certainty (100\% probability) differently than near-certainty (89\% or 90\%)
    \item Alternative models like Prospect Theory have been developed to account for such violations
\end{itemize}

\subsection{Why Independence Matters}

The independence axiom is crucial because:
\begin{enumerate}
    \item It's necessary (along with the other axioms) for preferences to be representable by expected utility
    \item It ensures logical consistency in decision-making
    \item Violations lead to predictable patterns of behavior that can be exploited
    \item It separates rational choice theory from descriptive theories of actual behavior
\end{enumerate}

\section{Summary}

\begin{itemize}
    \item The \textbf{independence axiom} states that preferences between lotteries should not change when both are mixed with a third lottery in the same proportion
    \item Sally's preferences show: $(0, 1, 0) \succ (0.01, 0.89, 0.10)$ and $(0.90, 0, 0.10) \succ (0.89, 0.11, 0)$
    \item These preferences lead to a \textbf{mathematical contradiction} when we try to represent them with expected utility
    \item The contradiction proves Sally violates the independence axiom
    \item Such violations are common in real-world decision-making and motivate alternative theories
\end{itemize}

\end{document}
